\documentclass[11pt]{article}
\usepackage[a4paper,margin=25mm]{geometry}
\usepackage[T1]{fontenc}
\usepackage{lmodern,amsmath,amssymb,amsthm,mathtools,microtype,booktabs,enumitem}
\usepackage[colorlinks=true,linkcolor=blue,urlcolor=blue,citecolor=blue]{hyperref}
\hypersetup{pdftitle={A fifth proof that zeta(3) is irrational},pdfauthor={Mathathon peer-review digest}}
\newcommand{\Z}{\mathbb Z}
\newcommand{\dd}{\,\mathrm d}
\DeclareMathOperator{\lcm}{lcm}
\newtheorem{theorem}{Theorem}
\newtheorem{lemma}{Lemma}
\newtheorem{proposition}[lemma]{Proposition}
\newtheorem{corollary}[lemma]{Corollary}
\setlist{itemsep=1pt,topsep=3pt}
\newcommand{\q}[1]{\medskip\noindent\textbf{#1}}
\title{A fifth proof that $\zeta(3)$ is irrational\\[2mm]
\large a maximum principle, and non-vanishing by a congruence}
\author{Mathathon peer-review digest\\
\small results of a short search by five AI agents, assembled and checked by Claude}
\date{14 November 2026}
\begin{document}
\maketitle

\noindent\textbf{What this note contains.} Two things came out of the search.
\begin{enumerate}
\item \textbf{An ending for the proof, new to this project} (Section~1). Let $I_n$ be the numbers of
the main document; they satisfy Ap\'ery's recurrence
$(n+1)^3u_{n+1}-(2n+1)(17n^2+17n+5)u_n+n^3u_{n-1}=0$. Once they are known to satisfy it and to be
bounded, their sign and size follow in half a page (Lemma~\ref{lem:max}) from a \emph{maximum
principle}: $I_n>0$ and $I_n<I_{n-1}/33$. No integral, no second sequence, no determinant identity,
no growth estimate. (These are the tools of the four endings in the main document, which is why we
call this one the fifth: Beukers' triple integral; the identity $n^3(a_nI_{n-1}-a_{n-1}I_n)=12$
together with the growth of the second solution $a_n$; the inequality
$I_n\le1/(n(n+1)(2n+1)a_n)$, stated there; and Ball's second series, used by Zudilin.) Together
with three facts quoted from the main document this is a complete proof using sums only.
\item \textbf{A way to see that $I_p\ne0$ for all large primes $p$, new to this project}
(Section~2), assuming $\zeta(3)=P/Q$: a congruence modulo $p$. It uses no inequality at all, so it
does not need the sign. With it we sketch a proof in which the sign of $I_n$ plays no role; one
classical Fourier integral is quoted without proof.
\end{enumerate}
\noindent\textbf{What it does not contain.} A proof that avoids Ap\'ery's numbers altogether. None of
the five agents found one (Section~3).

\smallskip\noindent\textbf{Status.} Found and checked by AI agents on 14 November 2026;
\emph{no human has read this}. Notation is that of the main document
\texttt{zeta3-proofs.pdf}, whose Sections 9 and 10 are quoted where indicated; the papers referred to
are Zudilin, arXiv:math/0202159, and Rajkumar, arXiv:1212.5881.

\smallskip\noindent\textbf{Relation to the literature.} The idea behind Section~1 is known; at
most the packaging is new. The details and references are at the end of Section~1. No search was
made for the congruence of Section~2.

\section{The proof by a maximum principle}

\q{The setting (recalled from the main document, Sections 2, 4 and 9).} Let
$D_n=\lcm(1,\dots,n)$ and
\[
r_n(t)=\frac{(t-1)(t-2)\cdots(t-n)}{t(t+1)\cdots(t+n)},\qquad R_n=r_n^2,\qquad
I_n=-\sum_{t=1}^\infty R_n'(t).
\]
\emph{Arithmetic.} Partial fractions give (main document, Section~9) $I_n=2(a_n\zeta(3)-b_n)$ with
$a_n=\sum_k\binom nk^2\binom{n+k}k^2$ an integer and $D_n^3b_n$ an integer. So if $\zeta(3)=P/Q$
with positive integers $P,Q$, then $\frac12Q\,D_n^3I_n=D_n^3(a_nP-b_nQ)$ is an integer.
\emph{What is needed:} $I_n\neq0$ and $|I_n|\le C\theta^n$ with $\theta<\frac1{27}$, because
$D_n<3^n$ for infinitely many $n$ (main document, Section~7).
\emph{The same numbers as integrals} (not needed for the proof; used only in three remarks below):
with $P_n(x)=\sum_{k=0}^n(-1)^k\binom nk\binom{n+k}kx^k$ and $K(x,y)=\frac{-\log(xy)}{1-xy}$ one has
$I_n=\int_0^1\!\!\int_0^1P_n(x)P_n(y)K(x,y)\dd x\dd y$ (main document, Sections 4 and 9).

\q{Step 1: the recurrence, written so that its structure shows.} The main document (Section~10)
finds, by solving five linear equations in six unknowns, the identity
\[
(n+1)^3R_{n+1}(t)-\beta_nR_n(t)+n^3R_{n-1}(t)=\sigma_n(t+1)R_n(t+1)-\sigma_n(t)R_n(t),
\]
with $\beta_n=(2n+1)(17n^2+17n+5)$ and $\sigma_n(t)=4(2n+1)(-2t^2+t+(2n+1)^2)$. Differentiate in
$t$ and sum over $t=1,\dots,T$: the right side becomes $(\sigma_nR_n)'(T+1)-(\sigma_nR_n)'(1)$. The
second term is $0$ because $R_n$ has a double zero at $t=1$ (this is where $n\ge1$ is used), and
the first is $O(T^{-2})$, because $\sigma_nR_n$ is a quotient of two polynomials of the same degree.
Letting $T\to\infty$ gives $(n+1)^3I_{n+1}-\beta_nI_n+n^3I_{n-1}=0$ for $n\ge1$; at $n=1$ this
involves $I_0=-\sum_t(t^{-2})'=2\zeta(3)$. Now observe that
\[
\beta_n=(n+1)^3+n^3+4(2n+1)^3 .
\]
(Indeed $(n+1)^3+n^3=(2n+1)(n^2+n+1)$ and $17n^2+17n+5=(n^2+n+1)+4(2n+1)^2$.) So the recurrence is
\begin{equation}\label{SL}
(n+1)^3\bigl(I_{n+1}-I_n\bigr)-n^3\bigl(I_n-I_{n-1}\bigr)=4(2n+1)^3\,I_n\qquad(n\ge1).
\end{equation}

\q{Why this form is the right one.} Read $I_{n+1}-I_n$ as a derivative. Then \eqref{SL} is the
discrete version of a differential equation $(p\,u')'=q\,u$ with $p>0$ and $q>0$. Such an equation
has a simple feature: \emph{a solution cannot have a positive maximum} (there $u'=0$ and
$u''\le0$, so $(pu')'=pu''\le0<qu$), and for the same reason it cannot have a negative minimum. So
a solution that is positive and starts to rise keeps rising, and one that has become negative
while falling keeps falling. If boundedness excludes both, a solution with $u(0)>0$ must stay
positive and decrease all the way. Lemma~\ref{lem:max} is this argument for sequences; there the
rise and the fall are shown to be unbounded.

\q{Step 2: $I_n$ is bounded.} For integers $t>n\ge1$ one has $0<r_n(t)<\frac1t$, and
$R_n'=2R_n\varphi_n$ where $\varphi_n(t)=\sum_{j=1}^n\frac1{t-j}-\sum_{j=0}^n\frac1{t+j}$ is a
difference of two positive numbers each at most $n$. Since $R_n'(t)=0$ for $t=1,\dots,n$, we get
$|I_n|\le2n\sum_{t>n}t^{-2}<2$ for $n\ge1$; and $I_0=2\zeta(3)<3$. So the sequence $(I_n)$ is
bounded. (This is the crude bound of the main document; nothing better is needed.)

\begin{lemma}[Maximum principle]\label{lem:max}
Let $(X_n)_{n\ge0}$ be a bounded real sequence with $X_0>0$ satisfying
\begin{equation}\label{star}
(n+1)^3(X_{n+1}-X_n)-n^3(X_n-X_{n-1})=4(2n+1)^3X_n\qquad(n\ge1).
\end{equation}
Then $X_0>X_1>X_2>\dots>0$, and $X_n<X_{n-1}/33$ for every $n\ge1$.
\end{lemma}
\begin{proof}
Put $d_n=X_{n+1}-X_n$, so that \eqref{star} reads $(n+1)^3d_n=n^3d_{n-1}+4(2n+1)^3X_n$.

\emph{Once up, always up.} Suppose $X_m>0$ and $d_m\ge0$ for some $m\ge0$. Then $X_{m+1}>0$, and
\eqref{star} at $n=m+1$ has a non-negative term and a positive term on the right, so $d_{m+1}>0$.
Repeating, $d_n>0$ and $X_n\ge X_{m+1}=:c>0$ for all $n>m$. Then, since $n^3d_{n-1}\ge0$,
$(n+1)^3d_n\ge4(2n+1)^3X_n\ge4(2n+1)^3c>4(n+1)^3c$ for all $n>m$, so $d_n>4c$ and $X_n\to\infty$.
Impossible.

\emph{Once non-positive, always down.} Suppose $X_m>0\ge X_{m+1}$ for some $m\ge0$. Then $d_m<0$, and
\eqref{star} at $n=m+1$ gives $(m+2)^3d_{m+1}=(m+1)^3d_m+4(2m+3)^3X_{m+1}<0$. So
$X_{m+2}<X_{m+1}\le0$. Repeating, $d_n<0$ and $X_n\le X_{m+2}=:-c<0$ for $n\ge m+2$, and then
$(n+1)^3d_n<-4(2n+1)^3c$, so $d_n<-4c$ and $X_n\to-\infty$. Impossible.

\emph{Conclusion.} If $X_m>0$, the first paragraph gives $X_{m+1}<X_m$ and the second gives
$X_{m+1}>0$. By induction from $X_0>0$, all $X_n$ are positive and the sequence decreases. Finally
write \eqref{star} as
\[
n^3(X_{n-1}-X_n)=(n+1)^3(X_n-X_{n+1})+4(2n+1)^3X_n .
\]
Both terms on the right are positive, and $(2n+1)^3>8n^3$, so $n^3(X_{n-1}-X_n)>32\,n^3X_n$, that
is, $X_{n-1}>33\,X_n$.
\end{proof}

\begin{theorem}
$0<I_n\le2\zeta(3)\,33^{-n}$ for all $n$, and $\zeta(3)$ is irrational.
\end{theorem}
\begin{proof}
By Steps 1 and 2 the sequence $I_n$ satisfies the hypotheses of Lemma~\ref{lem:max}, with
$I_0=2\zeta(3)>0$. If $\zeta(3)=P/Q$ with positive integers $P,Q$, then for the infinitely many
$n$ with $D_n<3^n$ the number
$\frac12QD_n^3I_n$ is a positive integer smaller than $Q\zeta(3)\,(27/33)^n$, which is less than $1$
for large $n$.
\end{proof}

\q{How good is the bound?} The lemma gives $I_{n-1}/I_n>1+4(2+\frac1n)^3$. For $n=1,2,3$ this says
$109$, $63.5$, $51.8$; the true ratios are $116.9$, $66.8$, $54.1$, and they decrease to
$17+12\sqrt2=33.97$. So the constant $33$ is nearly the true rate; the only estimate used is
$(2n+1)^3>8n^3$.

\q{What was needed, and what was not.} Needed: the arithmetic of $I_n$; one telescoping identity;
the bound $|I_n|<2$; the lemma; $D_n<3^n$ infinitely often. Not needed: any integral; the integers
$a_n$ beyond their integrality (in particular not their recurrence and not their growth); the
identity between two solutions used by Ap\'ery; a second series; an optimisation; the prime number
theorem.

\q{The same ending for other constants} (found by one of the agents; each checked exactly for
$n\le40$). \emph{For $\log2$.} With $P_n$ as above,
$y_n=\int_0^1\frac{P_n(x)}{1+x}\dd x=A_n\log2-B_n$ (where $A_n=P_n(-1)\in\Z$, $D_nB_n\in\Z$,
$y_0=\log2$) satisfies $(n+1)(y_{n+1}-y_n)-n(y_n-y_{n-1})=2(2n+1)y_n$ for $n\ge1$. (Divide the
three-term relation of the Legendre polynomials, $(n+1)P_{n+1}=(2n+1)(1-2x)P_n-nP_{n-1}$, by $1+x$
and integrate, using $\frac{1-2x}{1+x}=-2+\frac3{1+x}$ and $\int_0^1P_n=0$.) If $(y_n)$ is known to
be bounded, the proof of Lemma~\ref{lem:max} applies with $n$ in place of $n^3$ and $2(2n+1)$ in
place of $4(2n+1)^3$, and gives $0<y_n<y_{n-1}/(5+\frac2n)$; and $5>3$ is enough against $D_n$.
(Boundedness has to come from elsewhere. The formula
$y_n=\int_0^1\frac{x^n(1-x)^n}{(1+x)^{n+1}}\dd x$, formula (4) of the main document, gives it, but
it shows sign and size directly; so for $\log2$ the lemma is an alternative, not a shortcut.)
A streamlined, self-contained version of Section~1 is in the file \texttt{max-principle-proof.pdf}.

\emph{For $\zeta(2)$.} The integers $\alpha_n=\sum_k\binom nk^2\binom{n+k}k$ satisfy
$(n+1)^2u_{n+1}=(11n^2+11n+3)u_n+n^2u_{n-1}$; let $\beta'_n$ be the solution of the same recurrence
with $\beta'_0=0$, $\beta'_1=5$ (these are Ap\'ery's numbers for $\zeta(2)$, and
$D_n^2\beta'_n\in\Z$). This recurrence is not of the
type \eqref{star}: its small solution alternates in sign. But
$x_n=(-1)^n(\alpha_n\zeta(2)-\beta'_n)$ satisfies $n^2x_{n-1}=(n+1)^2x_{n+1}+(11n^2+11n+3)x_n$, and
\emph{if $x$ is bounded and $x_0>0$, then all $x_n>0$ and $x_n<x_{n-1}/11$}. Indeed, otherwise take
$M$ with $x_M>0\ge x_{M+1}$; then $s_k=(-1)^kx_{M+k}$ has $s_0>0$, $s_1\ge0$ and
$(n+1)^2s_{k+1}=(11n^2+11n+3)s_k+n^2s_{k-1}$ with $n=M+k$, so $s_2>0$, $s_{k+1}\ge\frac{11}4s_k$ and
$s_k\to\infty$. So all $x_n>0$, and then $n^2x_{n-1}>(11n^2+11n+3)x_n>11n^2x_n$. Finally $11>9$ is
enough against $D_n^2$. (Here too boundedness is an assumption.)

\q{A remark on the two-index table.} Let
$I_{n,m}=\int_0^1\!\!\int_0^1P_n(x)P_m(y)K(x,y)\dd x\dd y$, so that $I_{n,n}=I_n$ (according to the
agent these are the entries of the table in Rajkumar's paper). The same agent showed that, for each
fixed $m$,
\[
(n+1)^3(I_{n+1,m}-I_{n,m})+n^3(I_{n-1,m}-I_{n,m})=2(2n+1)\,m(m+1)\,I_{n,m}\qquad(n\ge1),
\]
again with a non-negative right-hand coefficient, and that sign and size of the diagonal can be
obtained from these row equations alone, with a little more work. The relation was verified exactly
for $m\le6$, $n\le10$ (it depends on the sign convention $P_n(0)=1$); we have not checked the
derivation of sign and size and do not use it.

\q{Relation to the literature} (short search, 14 November 2026; the passages cited were read). The
condition that makes Lemma~\ref{lem:max} work, that the middle coefficient of Ap\'ery's recurrence
exceeds the sum of the outer ones, was observed by A.~Mingarelli (arXiv:math/0608577), who calls the
recurrence disconjugate and identifies $a_n\zeta(3)-b_n$ as its recessive solution; he states that
this is a restatement of Ap\'ery's proof and takes the growth of $a_n$ from elsewhere. Mahboubi and
Sibut-Pinote (arXiv:1912.06611, Lemma~5.6) prove $33^n=O(a_n)$ and use $27<33$. Rajkumar
(arXiv:1212.5881) gets the rate from the recurrence for the linear forms alone, by the
Poincar\'e--Perron theorem. We did not find the elementary argument of Lemma~\ref{lem:max} applied
directly to $I_n$. So the idea is known; at most the packaging is new. No search was made for the
congruence of Section~2.

\section{Non-vanishing without any inequality}

\q{Question.} Every proof so far shows $I_n\ne0$ by showing $I_n>0$. Is the sign really needed?
In the language of sums the sign of $I_n$ is invisible (main document, Section~11: the terms of
$-\sum R_n'(t)$ have both signs), so a reason for non-vanishing that does not pass through the sign
is worth having.

\q{Idea.} Under the hypothesis $\zeta(3)=P/Q$ the number $N_n=D_n^3(a_nP-b_nQ)$ is an integer. An
integer is non-zero as soon as it is non-zero modulo some prime. Take $n=p$ prime. Then $D_p$ is
divisible by $p$ exactly once, so $D_p^3a_pP\equiv0\pmod p$, and everything depends on the part of
$b_p$ that has $p^3$ in its denominator. That part can be computed exactly.

\q{Congruences between fractions.} For rational numbers $x,y$ whose denominators in lowest terms
are not divisible by $p$, we write $x\equiv y\pmod p$ if the numerator of $x-y$ is divisible by
$p$; such fractions add and multiply like integers modulo $p$. A fraction ``has $p^e$ in its
denominator'' if its denominator in lowest terms is divisible by $p^e$.

\begin{proposition}\label{prop:cong}
Let $p\ge5$ be prime and $0\le m<p$. Then the denominator of $p^3\,b_{p+m}$ is prime to $p$, and
\[
p^3\,b_{p+m}\equiv6\,a_m\pmod p .\qquad\text{In particular}\quad p^3b_p\equiv6\pmod p .
\]
\end{proposition}
\begin{proof}
Let $n=p+m$, let $c_k=(-1)^k\binom nk\binom{n+k}k$ (the numbers $c_k$ of formula (9) of the main
document up to the common sign $(-1)^n$, which is irrelevant because only products occur), and
$H_k^{(j)}=\sum_{u=1}^ku^{-j}$. By Section~9 of the main document,
\[
b_n=\sum_{k}c_k^2H_k^{(3)}+\sum_k\sum_{j\ne k}\frac{c_kc_j}{j-k}\,H_k^{(2)} .
\]
As $n<2p$, the only multiple of $p$ among $1,\dots,n$ is $p$ itself, and the only multiples of $p$
among the differences $j-k$ are $\pm p$. So $p^3$ appears in a denominator in exactly two ways: in
$c_k^2H_k^{(3)}$ with $k\ge p$, through the part $c_k^2/p^3$; and in $\frac{c_kc_j}{j-k}H_k^{(2)}$
with $k\ge p$ (so that $H_k^{(2)}$ contains $\frac1{p^2}$) and $j=k-p$, through the part
$-c_kc_{k-p}/p^3$. All other parts have $p$ at most twice in the denominator and vanish modulo $p$
after multiplication by $p^3$. Hence the denominator of $p^3b_n$ is prime to $p$, and
\[
p^3b_n\equiv\sum_{k=p}^{n}c_k\,(c_k-c_{k-p})\pmod p .
\]
For $m=0$ the only term is $k=p$, with $c_0=1$ and $c_p=-\binom{2p}p\equiv-2$; the sum is
$(-2)(-3)=6$. For general $m$ write $k=p+l$ with $0\le l\le m$ and let
$\gamma_l=(-1)^l\binom ml\binom{m+l}l$, so that $a_m=\sum_l\gamma_l^2$. \emph{Lucas' theorem} says
$\binom{a_1p+a_0}{b_1p+b_0}\equiv\binom{a_1}{b_1}\binom{a_0}{b_0}\pmod p$ for $0\le a_0,b_0<p$. It
gives $\binom{p+m}{p+l}\equiv\binom{p+m}l\equiv\binom ml$. If $m+l<p$, it gives
$\binom{2p+m+l}{p+l}\equiv2\binom{m+l}l$ and $\binom{p+m+l}l\equiv\binom{m+l}l$. If $m+l\ge p$, put
$s=m+l-p$, so that $0\le s<l$; then $\binom{m+l}l=\binom{p+s}l$,
$\binom{2p+m+l}{p+l}=\binom{3p+s}{p+l}$ and $\binom{p+m+l}l=\binom{2p+s}l$ are all $\equiv0$, because
by Lucas each is congruent to a multiple of $\binom sl=0$; so the same two congruences hold. With
$(-1)^{p+l}=-(-1)^l$ we get $c_{p+l}\equiv-2\gamma_l$ and $c_l\equiv\gamma_l$, and therefore
$p^3b_n\equiv\sum_l(-2\gamma_l)(-3\gamma_l)=6\sum_l\gamma_l^2=6\,a_m$. (The proof uses only that $p$
is odd; $p\ge5$ is needed in Corollary~\ref{cor:nonvan}, where $6\not\equiv0$.)
\end{proof}

\begin{corollary}[Non-vanishing]\label{cor:nonvan}
Suppose $\zeta(3)=P/Q$. Let $p\ge5$ be a prime not dividing $Q$. Then the integer
$N_p=D_p^3(a_pP-b_pQ)$ is not divisible by $p$; in particular $I_p\ne0$.
\end{corollary}
\begin{proof}
$D_p=p\,D_{p-1}$ with $D_{p-1}$ prime to $p$. So $D_p^3a_pP\equiv0$ and
$D_p^3b_pQ=D_{p-1}^3\,(p^3b_p)\,Q\equiv6\,D_{p-1}^3Q\not\equiv0\pmod p$.
\end{proof}
\noindent No inequality and no recurrence was used, only the explicit formula for $b_n$. Showing
that an integer is non-zero by reducing it modulo a suitable prime is the device of Hilbert's 1893
proof that $e$ is transcendental (quoted from memory, sources not retrieved); here it is applied to
Ap\'ery's combinations.

\q{A proof in which the sign plays no role (sketch).} With Corollary~\ref{cor:nonvan}, a bound
$|I_p|\le C\theta^p$ with $\theta<e^{-3}$ along primes finishes the proof, once $D_p$ is controlled
along primes (see below). One of the agents found a bound that never looks at the sign. For every
$c\in\{\frac12,\frac32,\dots,n+\frac12\}$,
\begin{equation}\label{line}
I_n=\frac\pi2\int_{-\infty}^{\infty}\frac{r_n(c+iu)^2}{\cosh^2(\pi u)}\dd u,
\qquad\text{and hence}\qquad |I_n|\le\frac{\pi^2c}4\,r_n(c)^2 .
\end{equation}
\emph{The formula.} Let $c=k+\frac12$ with $0\le k\le n$. Since $R_n'(t)=0$ for $t=1,\dots,n$, we
may drop the first $k$ terms:
$I_n=-\sum_{t\ge k+1}R_n'(t)=\int_0^1\!\!\int_0^1P_n(x)P_n(y)\,(xy)^k\,\frac{-\log(xy)}{1-xy}\dd x\dd y$
(main document, Section~9). Insert the classical Fourier integral
$\frac{-\log z}{1-z}=\frac\pi2\int_{-\infty}^\infty\frac{z^{-1/2+iu}}{\cosh^2\pi u}\dd u$ ($z>0$)
with $z=xy$ and exchange the integrals (allowed because
$\iint|P_n(x)P_n(y)|(xy)^{k-1/2}<\infty$ and $\int\cosh^{-2}<\infty$). The inner integrals are
$\int_0^1P_n(x)x^{c-1+iu}\dd x=(-1)^nr_n(c+iu)$ by formula (8) of the main document, which holds for
complex exponents with positive real part because both sides are the same rational function. For
$c=\frac12$ the resulting bound is only $\frac{\pi^2}{2(2n+1)^2}$; discarding the first $k$ terms is
what makes it exponentially small. (For other real $c$ in $(0,n+1)$ the formula is false as it
stands; it holds with $\sin^2\pi(c+iu)$ in place of $\cosh^2\pi u$.)

\emph{The inequality.} We use the product
$\cosh\pi u=\prod_{h=\frac12,\frac32,\dots}(1+\frac{u^2}{h^2})$. First,
$\frac{|r_n(c+iu)|^2}{r_n(c)^2}=\prod_{j=1}^n\bigl(1+\frac{u^2}{(c-j)^2}\bigr)\big/
\prod_{j=0}^n\bigl(1+\frac{u^2}{(c+j)^2}\bigr)$. The numbers $|c-j|$, $1\le j\le n$, lie in
$\{\frac12,\frac32,\frac52,\dots\}$; a value $h$ occurs at most twice ($j=c\pm h$), and the value $h=c$ at most once (only
$j=2c$, since $j=0$ is not in the range). As $\cosh^2\pi u=\prod_h(1+\frac{u^2}{h^2})^2$ contains
every factor twice, the whole numerator cancels and at least one factor $1+\frac{u^2}{c^2}$ of
$\cosh^2$ is left; with the factor $j=0$ of the denominator this makes $(1+\frac{u^2}{c^2})^2$, and
all other remaining factors are $\ge1$. So
$\frac{|r_n(c+iu)|^2}{r_n(c)^2\cosh^2\pi u}\le(1+\frac{u^2}{c^2})^{-2}$, and
$\int_{-\infty}^\infty(1+\frac{u^2}{c^2})^{-2}\dd u=\frac{\pi c}2$.

\emph{The bound.} For odd $n$ and $c=\frac n2$ one finds $|r_n(\frac n2)|=\frac{2\,(n!!)^3}{n^2(3n)!!}$,
where $n!!=n(n-2)\cdots3\cdot1$. With $\rho_n=(n!!)^3/(3n)!!$ one has
$\frac{\rho_{n+2}}{\rho_n}=\frac{(n+2)^2}{27n^2+54n+24}\le\frac{n+2}{27n}$, hence
$\rho_n\le\frac n3\,27^{-(n-1)/2}$, and \eqref{line} gives
$|I_n|\le\frac{3\pi^2}{2n}\,27^{-n}$ for odd $n$.

\emph{The primes.} A \emph{proper prime power} is a number $q^a$ with $q$ prime and $a\ge2$. By
Lemma~6 of the main document (Section~7) there are infinitely many $n$ with $D_n<2.8^n$. For such
$n$ let $p$ be the least prime $\ge n$. The prime powers in $(n,p]$ are $p$ itself (if $p>n$) and
proper prime powers $q^a$; each of the latter multiplies the least common multiple by $q\le\sqrt p$,
and there are at most $\sqrt p\,\log_2p$ of them (for each $a=2,\dots,\lfloor\log_2p\rfloor$ at most
$p^{1/a}\le\sqrt p$ values of $q$; $\log_2$ is the logarithm to base $2$). So
$D_p\le D_n\cdot p\cdot(\sqrt p)^{\sqrt p\log_2p}\le2.8^p\,e^{\sqrt p(\ln p)^2}<2.9^p$ for large $p$.

\emph{Conclusion.} For the infinitely many primes $p$ with $D_p<2.9^p$, $p\ge5$ and $p\nmid Q$,
\[
1\le|N_p|=\tfrac Q2\,D_p^3\,|I_p|\le\frac{3\pi^2Q}{4p}\Bigl(\frac{2.9^3}{27}\Bigr)^p\to0,
\]
a contradiction. \emph{What is not proved here} is the Fourier integral (the Fourier transform of
$1/\cosh^2$), which is classical but not elementary; we verified it, \eqref{line} and the bound
numerically (Section~4). So this second proof is a sketch. Formula \eqref{line} is the classical
contour-integral representation of $I_n$ (Gutnik, Nesterenko) on a line where
$\sin^2\pi s=\cosh^2\pi u$. What differs here is the estimate (one inequality instead of the
saddle-point method) and the non-vanishing (the congruence instead of positivity). Nesterenko's
paper was not retrieved; the attribution is from memory and from Zudilin's survey
arXiv:math/0206176.

\section{What the search found, and did not find}

Five agents worked for about ten minutes each on different ideas. In brief:

\begin{center}\small
\begin{tabular}{@{}p{38mm}p{112mm}@{}}
\toprule
idea & outcome\\
\midrule
maximum principle & Section~1. Correct; extends to $\log2$ and $\zeta(2)$.\\
free search (one agent, 13 unrelated strategies) & Section~2: the congruence, and the sign-free
bound. Abandoned: constructions by Siegel's lemma, that is, with coefficients obtained from the
pigeonhole principle instead of explicitly (the non-vanishing step reduces to one congruence that
the lemma does not control); determinants of series remainders; continued-fraction criteria.\\
function theory in the half-plane & An explicit bound
$0<I_n\le\frac{\Gamma(2c)^4}{\Gamma(4c)}\bigl(\frac{\Gamma(n+1-c)}{\Gamma(n+1+c)}\bigr)^2$ for
$0<c<n+1$, within a factor $1.5$ of the truth at the best $c$ (the factor appears to tend to
$\sqrt2$). Its positivity step is Beukers' integral in another form, so it is not a new proof.\\
one formula showing everything & No new single positive series with visible $D_n^3$. A positive
double series with all three properties exists,
$I_n=\sum_{v\ge u\ge n+1}r_n(u)\binom{n+v-u}{n}\bigl(\frac{n!}{v(v+1)\cdots(v+n)}\bigr)^2$, but it is
Sorokin's integral expanded term by term (a known construction), and its equality with $I_n$ was
only checked numerically for $n\le5$. Ball's series (main document, Section~11) equals $-R_n'$ plus
a telescoping term, exactly, for $n\le5$, but the telescoping function has fourth-order poles and no
closed form, so this does not explain $D_n^3$.\\
is our combination the only one? & No. See below.\\
\bottomrule
\end{tabular}
\end{center}

\q{The classical combination is not the smallest one.} The fifth agent considered all sums
$\sum_{t\ge1}F(t)=A\zeta(3)-B$ with
$F(t)=\sum_{k=0}^n\bigl[\frac{A_k}{(t+k)^3}+\frac{B_k}{(t+k)^2}+\frac{C_k}{t+k}\bigr]$, where $A_k$,
$D_nB_k$, $D_n^2C_k$ are integers and $\sum B_k=\sum C_k=0$ (the class containing $-R_n'$; then
$A\in\Z$ and $D_n^3B\in\Z$). Every pair $(A,B)\in\Z\times\frac1{D_n^3}\Z$ occurs, already with
$B_k=C_k=0$: then $A=\sum_kA_k$ and $B=\sum_{j=1}^nT_j/j^3$ with arbitrary integers
$T_j=\sum_{k\ge j}A_k$, and the numbers $D_n^3/j^3$ ($1\le j\le n$) have no common factor, so these
sums fill $\frac1{D_n^3}\Z$. (So this is a fact about the tails $\sum_{t>k}t^{-3}$ of the series, and
$I_n$ is one point of a set that contains everything.) In particular far smaller combinations than
$I_n$ exist: for $n=10$ and $|A|\le2a_{10}$ the smallest value is $1.4\cdot10^{-24}$, at
$A=668\,163\,577\,273$, against $I_{10}=2.2\cdot10^{-17}$. (The restriction on $A$ is only a
normalisation: the criterion looks at $D_n^3|A\zeta(3)-B|$, not at $A$.) These are the
continued-fraction approximations of $D_n^3\zeta(3)$. A single such value, certified numerically,
proves only that a denominator of $\zeta(3)$ exceeds $1/(D_n^3\cdot\text{value})$, here
$4\cdot10^{13}$; irrationality needs $D_n^3\cdot|\text{value}|\to0$ along an infinite family, proved
for all $n$ at once. \emph{What singles out Ap\'ery's combination is not that it is small but that
its smallness can be proved.} So the objects that the main document calls ``forced'' are not singled
out by size; what distinguishes them is that their size can be proved. This is consistent with the
search for a proof with different numbers finding nothing: there is no shortage of small
combinations, only of provable ones.

\section{Checks}
All by AI agents; no human.

\emph{First round.} The lemma: proof read line by line by two agents and by the orchestrator; the
conclusions $I_n>0$, $I_n<I_{n-1}/33$ verified with exact $a_n,b_n$ and $150$-digit arithmetic for
$n\le40$; the identity for $\beta_n$ symbolically. The congruence of Proposition~\ref{prop:cong}:
derived by hand (case $m=0$ by an agent, general $m$ by the orchestrator) and verified exactly in
$90$ cases ($p\le23$, $p+m\le40$), and for $m=0$ for all primes up to $43$. Formula \eqref{line}:
verified to $12$ digits for $n=1,3,5,7$ and three values of $c$.

\emph{Referee round.} A fresh agent read the note twice, for correctness and as the intended reader
($31$ points, all applied; no mathematical claim had to be withdrawn), and recomputed
independently: the recurrence exactly for $a_n$ and $b_n$ ($n\le44$), and $I_n>0$,
$I_n<I_{n-1}/33$ with $200$-digit arithmetic for $n\le45$ (about $125$ digits cancel at $n=40$, so
the $150$ digits of the first round were just enough); the telescoping identity symbolically in $n$
and $t$, and its boundary term for $n\le4$; Lemma~\ref{lem:max} by experiment (changing $I_1$ by
$\pm10^{-30}$ sends the solution to $\pm3\cdot10^{27}$ at $n=40$); Proposition~\ref{prop:cong} and
both Lucas congruences exactly in $127$ cases (all primes $5\le p\le31$ with $p+m\le45$; $46$ of
them have $m+l\ge p$ for some $l$) and for $p=2,3$; formula \eqref{line} for all $n\le8$ and all
admissible $c$ (relative error below $10^{-17}$), its failure for other $c$, the inequality at
$4620$ points, the Fourier integral at five values of $z$, and the bound
$\frac{3\pi^2}{2n}27^{-n}$ for odd $n\le301$; the statement on primes for all $n<298\,000$ (the
factorisation of $D_p/D_n$, the count, and $D_p<2.9^p$); the $\log2$ and $\zeta(2)$ statements
exactly for $n\le41$; the two-index relation for $m\le6$, $n\le10$; in Section~3 the half-plane
bound for $n\le45$, the value $1.4\cdot10^{-24}$, and the double series of the table numerically for
$n\le3$.

\emph{Not proved here:} the Fourier integral behind \eqref{line}; boundedness of the $\log2$ and
$\zeta(2)$ forms. \emph{Not verified:} the statements about Rajkumar's table and about Sorokin's
integral and Ball's series in the table of Section~3; the outcomes of the abandoned strategies; the
attributions to Hilbert, Gutnik and Nesterenko, which are from memory. For the paragraph on the
literature the referee confirmed only that arXiv:math/0608577 and arXiv:1912.06611 exist and are on
the stated subjects.

\end{document}
